\begin{abstract}
Computer-use agents (CUAs) are evaluated on whether they eventually reach the correct end state, as judged by rubrics or deterministic state checks. Users judge them differently: an agent that needs three minutes and forty actions to change a setting that a person changes in five seconds is correct, but nobody would wait for it. We introduce \bench, a benchmark that measures when CUAs frustrate the people they act for. It contains 75 human-written tasks that people complete in seconds but that are prone to agent failure, such as dragging an item, choosing from a hover menu, picking a date or dismissing a dialog. For every task, human annotators record the effort it takes them and the tolerance they would extend to an agent: how long, and through how many unproductive actions, they would watch it before stopping it and doing the task themselves. A run succeeds only if it is correct and stays within this tolerance, so failing, looping, stalling and overrunning the human budget all count as frustration. We also define eight frustration patterns, with detectors, that explain why a run frustrates. Across \tbd{M} frontier models, \tbd{the best model is correct on XX\% of runs but succeeds within human tolerance on only XX\%}; \tbd{correct runs take a median X.X$\times$ the human time}, and \tbd{looping} appears in \tbd{XX\%} of frustrated runs. Ranking models by tolerated success instead of correctness \tbd{reorders N of the M models}. The gap between agents and people is therefore larger than pass rates suggest. We release the tasks, human labels and evaluation harness.
\end{abstract}
